% !PRINTED none
\weberpage{none}
\begin{center}
\gesperrt{Fifth Section.}

\gesperrt{Constitution of the General Groups.}
\end{center}

\rule{4cm}{0.4pt}

\begin{center}
\S.\,27.

\gesperrt{Formation of Groups according to Cayley.}
\end{center}

The general definition of the group, which was given in \S.\,1, still leaves much about the nature of this concept in the dark, and the various individual groups too, which we have come to know in the course of our considerations, give only hints of general laws, and show that the concept of a group contains nothing contradictory in itself. More is contained in the definition of the group than appears at first sight, and the number of possible groups that can be composed from a given number of elements is a very limited one. The general laws which prevail here have as yet been recognised only to the smallest part, so that every new special group, especially one of smaller number of terms, offers a new interest and invites detailed study.

\gesperrt{What groups are at all possible among a given number of elements, that is, for a given degree?} That is the general question at issue, from whose complete solution, however, we are still far removed. \gesperrt{Cayley} first took up this problem for the lowest degrees\footnote{On the theory of groups, as depending on the symbolic equation $\Theta^n = 1$. Philosophical Magazine, Vol.\ II, 1854. (\gesperrt{Cayley's} mathematical papers, Vol.\ II, 125.) American Journal of mathematics, Vol.\ I.}.
